%=========== ΛΥΣΗ - 2ο ΘΕΜΑ Δ ===========
\begin{thema}{Δ}
Η συνάρτηση $f$ είναι δύο φορές παραγωγίσιμη στο $(-1,+\infty)$ ως πράξεις παραγωγίσιμων συναρτήσεων και για κάθε $x>-1$ είναι
\[
f'(x)=2e^{2x}-\frac{1}{x+1}
\]
και
\[
f''(x)=4e^{2x}+\frac{1}{(x+1)^2}.
\]
\begin{erwthma}
\item Για κάθε $x>-1$ είναι
\[
e^{2x}>0
\quad\text{και}\quad
\frac{1}{(x+1)^2}>0,
\]
οπότε
\[
f''(x)=4e^{2x}+\frac{1}{(x+1)^2}>0.
\]
Επομένως η $f$ είναι \textbf{αυστηρά κυρτή} στο διάστημα $(-1,+\infty)$.

\item Επειδή $f''(x)>0$ για κάθε $x>-1$, η $f'$ είναι γνησίως αύξουσα στο $(-1,+\infty)$. Επίσης
\[
f'\left(-\frac{1}{2}\right)
=2e^{-1}-\frac{1}{\frac{1}{2}}
=\frac{2}{e}-2<0
\]
και
\[
f'(0)=2e^0-1=1>0.
\]
Η $f'$ είναι συνεχής στο $\left[-\frac{1}{2},0\right]$. Επομένως, σύμφωνα με το Θεώρημα Ενδιάμεσων Τιμών, υπάρχει $x_0\in\left(-\frac{1}{2},0\right)$ τέτοιο ώστε
\[
f'(x_0)=0.
\]
Η ρίζα αυτή είναι μοναδική, διότι η $f'$ είναι γνησίως αύξουσα στο $(-1,+\infty)$.

\item Επειδή η $f'$ είναι γνησίως αύξουσα και μηδενίζεται στο $x_0$, έχουμε:
\begin{itemize}
\item αν $-1<x<x_0$, τότε $f'(x)<f'(x_0)=0$,
\item αν $x_0<x<+\infty$, τότε $f'(x)>f'(x_0)=0$.
\end{itemize}
Στον παρακάτω πίνακα βλέπουμε το πρόσημο της $f'$ και τη μονοτονία της $f$:
\begin{center}
\begin{tikzpicture}
\tkzTabInit[color,lgt=2,espcl=1.8,colorC = maincolor!50!white,colorL = maincolor!20!white,colorV = maincolor!50!white]%
{$x$ / .8,$f'(x)$ /0.8,$f(x)$/0.8}%
{$-1$,$x_0$,$+\infty$}%
\tkzTabLine{d, -, z, +, }
\tkzTabLine{d, \searrow, t, \nearrow, }
\end{tikzpicture}
\end{center}

Άρα η $f$ είναι \textbf{γνησίως φθίνουσα} στο διάστημα $(-1,x_0]$ και \textbf{γνησίως αύξουσα} στο διάστημα $[x_0,+\infty)$. Επομένως παρουσιάζει ολικό ελάχιστο στο $x_0$, με τιμή
\[
f(x_0)=e^{2x_0}-\ln(x_0+1)-1.
\]
Συνεπώς, για κάθε $x>-1$ ισχύει
\[
{f(x)\geq f(x_0)},
\]
με την ισότητα να ισχύει μόνο για $x=x_0$.

\item Έστω $x>0$. Εφαρμόζοντας το Θεώρημα Μέσης Τιμής για την $F$ στα διαστήματα $[x,2x]$ και $[2x,3x]$, υπάρχουν $\xi_1\in(x,2x)$ και $\xi_2\in(2x,3x)$ τέτοια ώστε
\[
F(2x)-F(x)=F'(\xi_1)x=xf(\xi_1)
\]
και
\[
F(3x)-F(2x)=F'(\xi_2)x=xf(\xi_2).
\]
Επειδή
\[
0<x<\xi_1<2x<\xi_2<3x
\]
και $x_0<0$, τα $\xi_1,\xi_2$ ανήκουν στο διάστημα $(x_0,+\infty)$, στο οποίο η $f$ είναι γνησίως αύξουσα. Άρα
\[
\xi_1<\xi_2\Rightarrow f(\xi_1)<f(\xi_2).
\]
Επειδή $x>0$, προκύπτει
\[
xf(\xi_1)<xf(\xi_2),
\]
δηλαδή
\[
F(2x)-F(x)<F(3x)-F(2x).
\]
Επομένως
\[
{F(3x)+F(x)>2F(2x)}.
\]
\end{erwthma}
\end{thema}
%=========== ΤΕΛΟΣ ΛΥΣΗΣ - 2ο ΘΕΜΑ Δ ===========
